
\section{Conclusion}
\label{sec:conclusion}
In this paper, we investigated the limitations of existing Acoustic Side-Channel Attack (ASCA) methods in controlled noisy noisy scenarios and demonstrated the potential of transformer architectures to address these challenges. By fine-tuning a CoAtNet model and introducing Vision Transformers (VTs), we established a new state-of-the-art benchmark in keystroke classification accuracy, surpassing traditional CNNs in clean acoustic conditions. Crucially, we showed that Large Language Models (LLMs) significantly mitigate errors caused by realistic noise, improving the reliability and practicality of ASCAs.
Furthermore, we found that lightweight LLMs fine-tuned with efficient techniques such as LoRA and QLoRA can achieve correction performance comparable to larger models, making them suitable even for resource-constrained attacks. Our results underline the importance of rigorously evaluating noise robustness in future ASCA research, highlighting transformer-based context-aware methods as essential for effective acoustic side-channel text recovery.
Finally, to enhance transparency and reproducibility, we will release our trained models, fine-tuning configurations, experimental pipeline code, and augmented acoustic-keyboard datasets with realistic noise scenarios as open-source resources. This aims to facilitate future research, encouraging the security community and practitioners to build on these findings toward comprehensive, noise-resilient ASCA solutions.
%In this paper, we systematically investigated the limitations of existing Acoustic Side-Channel Attack (ASCA) methods in realistic noisy scenarios and demonstrated the transformative potential of using transformer architectures to overcome these shortcomings. By carefully tuning an existing CoAtNet model and introducing Vision Transformers (VTs), we established a new state-of-the-art keystroke classification accuracy benchmark surpassing previous CNN models in clean acoustic scenarios. More importantly, our evaluation proved that Large Language Models (LLMs)—leveraging their powerful contextual reasoning capabilities—can effectively mitigate errors caused by realistic acoustic noise, greatly enhancing the reliability and practical feasibility of ASCAs.

%Moreover, we demonstrated that fine-tuned lightweight LLMs (through efficient methods such as LoRA and QLoRA) can achieve comparable correction performance to much larger language models, thus making such techniques viable even for deployment in resource-limited attack scenarios. Our results emphasize that future ASCA research must rigorously evaluate noise robustness, and transformer-driven context-aware approaches should be viewed as not merely beneficial but essential steps toward truly effective acoustic side-channel text recovery.

%Finally, to promote transparency, reproducibility, and future research opportunities, we commit to releasing our trained models, fine-tuning configurations, experimental pipeline code, and enhanced acoustic-keyboard datasets with noise scenarios as open-source resources in the final paper version. Our goal is to streamline subsequent advancements in acoustic side-channel research, encouraging both the security community and practitioners to build upon these foundational findings towards comprehensive and noise-resilient ASCA solutions.

\noindent \textbf{Reproducibility.} Our source code is available at \url{https://github.com/Botacin-s-Lab/EchoCrypt}, and our fine-tuned model weights are available at \url{https://huggingface.co/seyyedaliayati/zoom_model} and \url{https://huggingface.co/seyyedaliayati/phone_model}.

\noindent \textbf{Acknowledgments.} We would like to thank the NSF for the support via the CNS 2327427 grant.